\chapter*{Overview}
\addcontentsline{toc}{chapter}{Overview}
\markboth{Overview}{}

Poraqu\^e predicts the electronic structure of a crystal from its geometry
alone. Given a \file{POSCAR}, an \file{INCAR} and a \file{POTCAR}, it produces
the valence charge density and the kinetic energy density --- with no
wavefunctions and no self-consistency cycle.

It does this by chaining two learned operators:
\[
  \{\text{\file{POSCAR}},\ \text{\file{INCAR}},\ \text{\file{POTCAR}}\}
  \;\xrightarrow{\ \text{analytic}\ }\;
  \text{\file{EXTCAR}}
  \;\xrightarrow{\ \text{Model 1}\ }\;
  \text{\file{CHGCAR}}
  \;\xrightarrow{\ \text{Model 2}\ }\;
  \text{\file{TAUCAR}}
\]
The first step is closed-form; only the two field-to-field maps are learned.
Every output is written in \file{CHGCAR} format.

\section*{Who this guide is for}

This is the practical guide: how to lay out data, train, predict, and read the
results. It assumes you can run a plane-wave DFT calculation and want to use
Poraqu\^e as a tool.

For the physics and the architecture --- why a Fourier neural operator, what
the models correspond to in density-functional theory, how the external
potential is reconstructed from the pseudopotential tables --- see the
companion \emph{Technical Guide} in \file{latex/technical\_guide/}.

\section*{What you need}

\begin{center}
\begin{tabular}{@{}ll@{}}
\toprule
Requirement & Note \\
\midrule
Python 3.11 or newer & see Chapter~\ref{ch:install} \\
A set of DFT calculations & \file{CHGCAR} and \file{TAUCAR} per structure \\
An accelerator (optional) & CUDA or Apple Metal; CPU works \\
\opt{latexmk} (optional) & only for the automatic PDF reports \\
\bottomrule
\end{tabular}
\end{center}

\begin{pwarn}
Poraqu\^e learns from the data you give it. The published results were obtained
on seventeen gold supercells, and they measure interpolation between geometries of
one element --- plus, weakly, extrapolation across cell size, where the error
roughly doubles. Nothing in this guide should be read as evidence that a model
trained on one chemistry transfers to another --- validate on your own held-out
structures before trusting a prediction.
\end{pwarn}
